\section{Fysikingenjörens verktyg}

\subsection{Fel, mätosäkerhet}

\subsubsection{Noggranhet och precision}
\begin{figure}[!ht]
	\centering
	\includegraphics[width=10cm]{res/acc_vs_pres.jpg}
\end{figure}
Precision innebär att mät

\subsubsection{Felkällor och klassificering av fel}
Två sätt att klasisificera felkällor är enligt grova fel och systematiska fel.
Dessa innebär:
\begin{itemize}
	\item Grova fel \begin{itemize}
		\item Felaktigt handhavande
		\item Slarv
		\item Räknefel
	\end{itemize}
	\item Systematiska fel \begin{itemize}
		\item Instrumentfel
		\item Avläsning
		\item Yttre faktorer
	\end{itemize}
\end{itemize}

\subsubsection{Mätseries, medelvärden och spridningsmått}
Standardavvikelsen $\sigma$ är en kvantifisering av spridningsmåttet, och ges av
variansen:
$$
	\sigma^2=\frac{1}{n}\sum_{i=1}^n(x_i-\bar{x})^2
$$
Ett bra mått på mätnoggrannheten är standardavvikelsen i medelvärdet som ges av:
\begin{align*}
	\sigma_m&=\frac{1}{\sqrt{n(n-1)}}\sqrt{\sum_{i=1}^{n}(x_i-\bar{x})^2}\\
	&=\frac{1}{\sqrt{n(n-1)}}\sqrt{n\sigma^2}\\
	&=\sqrt{\frac{n\sigma^2}{n(n-1)}}\\
	&=\frac{\sigma}{\sqrt{n-1}}
\end{align*}

När man skriver en raport är det viktigt att berätta:
\begin{itemize}
	\item Metoden för beräkningen av felet onoggranhet
	\item \hl{Något mer}
\end{itemize}

\newpage

\subsubsection{Felfortplantning}
Felfortplantning är en teknik som behövs för att ange mätosäkerhet för storhet som
ej mäts direkt, utan fås genom funktionssamband. Vi vill få reda på storhet $y$
vilket beror på $x$ genom sambandet $y=f(x)$. Mätning ger $x=x_0\pm\Delta x$.

För att skatta $\Delta y$. Genom att linjarisera $f(x)$ kring $x$ får vi:
\begin{align*}
	y&=f(x_0\pm\Delta x)=f(x_0)\pm\qty|\dv{f}{x}|_{x_0}\Delta x\\
	&=y_0\pm\Delta y\\
	\implies y&\left\{\begin{matrix}
		y_0=f(x_0)\\
		\Delta y=\left|f'(x_0)\right|\Delta x
	\end{matrix}\right.
\end{align*}

Om vi å andra sidan har ett samband av flera variabler, t.ex. $y$ beror på $n$
oberoende variabler $y=f(x_1,x_2,\dots,x_n)$. Vi har då:
\begin{align*}
	x_1&=x_{10}\pm\Delta x_1\\
	x_2&=x_{20}\pm\Delta x_2\\
	&\qquad\vdots\\
	x_n&=x_{n0}\pm\Delta x_n
\end{align*}
Detta ger att:
$$
	y=f(x_{10},x_{20},\dots,x_{n0})\pm\sum_{k=1}^{n}\qty|\pdv{f}{x_k}|\Delta x_k
$$

\begin{ber}
	\textbf{Exempel på beräkning enligt felfortplantningsformeln}

	För en matematisk pendel gäller det för pendeltiden att:
	\begin{align*}
		T&=2\pi\sqrt{\frac{l}{g}}\\
		\implies T^2&=4\pi^2\frac{l}{g}\iff g=\frac{4\pi^2l}{T^2}
	\end{align*}
	Mätning ger:
	\begin{align*}
		l&=l_0\pm\Delta l\\
		T&=T_0\pm\Delta T
	\end{align*}
	Felfortplaningsformeln ger:
	\begin{align*}
		g(l,T)&\approx g(l_0,T_0)\pm\qty(\qty|\pdv{g}{l}|_{(l_0,T_0)}\Delta l+\qty|\pdv{g}{T}|_{(l_0,T_0)}\Delta T)\\
		&=g_0\pm\qty(\frac{4\pi^2}{T_0^2}\Delta l+\frac{8\pi^2l_0}{T_0^3}\Delta T)\\
		&=g_0\pm\qty(\frac{g_0}{l_0}\Delta l+\frac{2g_0}{T_0}\Delta T)\\
		\implies\Delta g&=g\qty(\frac{\Delta l}{l_0}+\frac{2\Delta T}{T_0})
	\end{align*}
	Relativ onoggranhet: $g$
\end{ber}

\subsection{Laminärt flöde}
\begin{itemize}
	\item Laminärt flöde innebär att strömlinjer ej korsar varandra.
	\item Turbulent flöde innebär motsatsen till laminärt flöde.
\end{itemize}
Laminärt flöde fås för bred långsamtflytande fluid.